\begin{proposition}[P4R.6]\hypertarget{tag-02YI}{}\label{night:P4R.6}\label{stmt:P4R.6}
\StatementRecord{night:P4R.6}{proved}{02YI}{\SeriesRef{night:P4.12}}

Let $S$ be an irreducible singular support in
Part~4.2, let $\nu:N\to S$ be its normalization,
and let $\rho:T\to N\to S$ use the minimal
resolution of $N$. There is an effective integral
divisor $D$ on $T$ such that
\[
 \omega_T(D)\simeq\rho^*\omega_S,
 \qquad \operatorname{Supp}D
       =\rho^{-1}(\operatorname{Sing}S)_{\rm red}.
\]
In particular $K_T+D$ is torsion, by P4.12.
The strict-transform conductor has its usual
positive integral coefficients. Every connected
exceptional configuration over a singular point
of $N$ has strictly positive coefficients in $D$.
Every such configuration meets the strict-transform
conductor.
\end{proposition}
\begin{proof}

\begin{enumerate}[label={$\langle1\rangle$\arabic*.},leftmargin=*]
\item The reduced support is a Cartier hypersurface in
the smooth threefold, so $\omega_S$ is invertible.
Choose a local generator $\eta$ of this line
bundle. On $S_{\rm reg}$ it is an ordinary
holomorphic two-form; its pullback is a
meromorphic two-form on $T$. Its divisorial
orders are integers. Two generators differ by
a holomorphic unit on $S$, whose pullback is
still a unit. Consequently their negative
divisors define one integral divisor $D$ and
the displayed line-bundle isomorphism.
Finite duality, in the form
\[
 \nu_*\omega_N
  =\mathcal Hom_{\mathcal O_S}
       (\nu_*\mathcal O_N,\omega_S),
\]
identifies the generic strict-transform pole
ideal with the conductor. This is the same
normalization adjunction used in P4.12. Thus
the coefficients away from the exceptional
curves are the positive integral conductor
exponents; away from conductor and exceptional
curves the pullback has no zero or pole.
No assertion that $K_N$ alone is
$\mathbb Q$-Cartier is needed.

\item We prove positivity along an exceptional
configuration. Its intersection matrix $M$
is negative definite, and we give the local
comparison establishing the matrix property
used here. Choose a nonzero local function
on the normal surface $N$ vanishing at the
image point. Its pullback divisor is
$H+\sum_i n_iE_i$, with $n_i>0$ and $H$
effective without exceptional components.
The strict transform meets the exceptional
configuration, because a curve through the
image point has closure meeting its proper
inverse image. The divisor is principal
near the configuration, so
\[
 Mn=-h,\qquad h_i=H\cdot E_i\ge0,
 \qquad h\ne0.
\]
The off-diagonal entries of $M$ are
nonnegative. On a connected configuration,
writing a real vector as $v_i=n_iw_i$
gives the identity
\[
 v^tMv
 =-\sum_{i<j}M_{ij}n_in_j(w_i-w_j)^2
   -\sum_i n_ih_iw_i^2.
\]
It is strictly negative for every nonzero
$v$: equality forces all $w_i$ equal along
the connected graph and then forces that
common value to be zero. The same maximum
principle, applied to $M(n_iw_i)=-b$ with
$b_i\ge0$, shows that $-M^{-1}$ has
nonnegative entries, strictly positive on
each connected block. This also follows
directly by taking a negative minimum of
the $w_i$ in the displayed weighted
Laplacian identity.

\item Write $D=D_h+\sum_i d_iE_i$, with $D_h$
the strict conductor. Minimality of the
resolution gives
\[
 K_T\cdot E_i=2p_a(E_i)-2-E_i^2\ge0.
\]
Indeed the only exception for a negative
curve of arithmetic genus zero would be
a smooth rational exceptional curve of
self-intersection $-1$, which is absent
from a minimal resolution. Also
$D_h\cdot E_i\ge0$. Intersecting
$K_T+D=\rho^*K_S$ with the contracted
curves gives
\[
 Md=-b,
 \qquad b_i=(K_T+D_h)\cdot E_i\ge0.
\]
Hence $d_i\ge0$. A zero coefficient on
a connected block forces $b=0$ there.
Every curve in that block is then a
smooth rational $(-2)$-curve, and the
block is disjoint from the strict
conductor.

\item In fact, every exceptional block disjoint from
the strict conductor is forbidden, independently
of its coefficients. Finite duality identifies
the conductor on the normal $N$, up to the
invertible pullback of $\omega_S$, with the
rank-one reflexive sheaf $\omega_N$. Thus its
support has no isolated point. Any conductor
curve through the image point has strict
transform meeting the exceptional block:
properness makes the closure of its punctured
inverse image meet the fibre. Therefore the
image of a block disjoint from the strict
conductor is away from the conductor, and
$\nu$ is an isomorphism near it. The
corresponding point of $S$ is a normal isolated
hypersurface singularity. The local
Milnor-fibre argument of P4.12 applies
at this point itself. A local reduced
equation $g$ has an isolated critical
point and positive Milnor number. A
Morse perturbation supplies a vanishing
two-sphere of square $-2$ in its Milnor
fibre. The original map is a unit times
$g^m$; taking a local root of the unit,
each of its nearby local fibres is a
disjoint union of copies of a Milnor
fibre of $g$. Such a copy embeds into
the nearby four-torus. Every map
$S^2\to T^4$ is null-homotopic, so that
embedded sphere would have zero
homology class and square zero, a
contradiction. Thus every exceptional
block has positive pole coefficients.
This proves both effectiveness and the
asserted support equality.

\end{enumerate}
\end{proof}